\appendix

\chapter{REVISÃO SISTEMÁTICA DA LITERATURA}
\label{apendice:rsl}

Apresente a revisão sistemática da literatura.

\chapter{LOCALIZAÇÃO ESPACIAL DOS PONTOS DE TRATAMENTO}
\label{apendice:localizacao_tratamento}

Apresente mapas e figuras dos pontos de tratamento.

\chapter{LOCALIZAÇÃO ESPACIAL DOS PONTOS DE CONTROLE}
\label{apendice:localizacao_controle}

Apresente mapas e figuras dos pontos de controle.

\chapter{TOTAL DE ACIDENTES EM CADA CONJUNTO DE PONTOS}
\label{apendice:total_acidentes}

Apresente tabelas com o total de acidentes.

\chapter{RELATÓRIO DO TESTE T}
\label{apendice:teste_t}

Apresente os resultados detalhados do teste estatístico.

\chapter{BATERIA DE TESTES: TRATAMENTO DE NOMES COMPLEXOS}
\label{apendice:teste_nomes}

Este apêndice documenta a bateria de testes automatizados para validação do motor de formatação de nomes de autores, extração de iniciais e tratamento de partículas conectivas (nacionais e estrangeiras) conforme implementado no pacote de formatação.

\begingroup
\newcommand{\testanome}[2]{%
  \begingroup
    \ultimosobrenomeautor{#1}%
    \prenomesautor{#2}%
    \noindent\textbf{Entrada:} \texttt{\textbackslash prenomesautor\{#2\}} e \texttt{\textbackslash ultimosobrenomeautor\{#1\}} \\
    \textbf{Cessão de Direitos (Completo):} \nomecompletoautor \\
    \textbf{Ficha Catalográfica:} \autorreferenciacatalografica \\
    \textbf{Referência da Dissertação:} \autorreferenciabibliografica \par
    \vspace{0.15cm}\hrule\vspace{0.3cm}
  \endgroup
}

\section*{Casos de Teste Avaliados}

% Teste 1: Padrão brasileiro com Grau de Parentesco
\noindent\textbf{Teste 1: Padrão brasileiro com Grau de Parentesco}\par\vspace{0.1cm}
\testanome{Silva Júnior}{José Augusto da}

% Teste 2: Tamanho extremo com múltiplas partículas
\noindent\textbf{Teste 2: Tamanho extremo com múltiplas partículas}\par\vspace{0.1cm}
\testanome{Bragança}{Pedro de Alcântara Francisco Antônio João Carlos Xavier de Paula Miguel Rafael Joaquim José Gonzaga Pascoal Cipriano Serafim de}

% Teste 3: Hífen (O motor extrai o 'J' de Jean-Jacques corretamente)
\noindent\textbf{Teste 3: Prenome composto com hífen}\par\vspace{0.1cm}
\testanome{Rousseau}{Jean-Jacques}

% Teste 4: Partícula e Sobrenome duplo hífen
\noindent\textbf{Teste 4: Partícula e sobrenome duplo composto}\par\vspace{0.1cm}
\testanome{García-Márquez}{Gabriel José de la Concordia}

% Teste 5: Partículas Germânicas e Holandesas (van, der, von)
\noindent\textbf{Teste 5: Partículas germânicas e holandesas (\textit{van}, \textit{der}, \textit{von})}\par\vspace{0.1cm}
\testanome{Gassen}{Win van der}
\testanome{Helmholtz}{Hermann von}

% Teste 6: Partículas Italianas e Espanholas (di, y)
\noindent\textbf{Teste 6: Partículas italianas e espanholas (\textit{di}, \textit{y})}\par\vspace{0.1cm}
\testanome{Vinci}{Leonardo di ser Piero}
\testanome{Cervantes}{Miguel de Cervantes y Saavedra}

\endgroup
